\section{Knowledge acquisition: parsing, ingestion and index verification}
\label{sec:acquisition}

Acquisition determines what later retrieval can support, and it is where
content comes to govern placement. QDCVR submits PDFs,
Office files and scans to MinerU through an asynchronous parsing service
\citep{mineru}; parse tasks return identifiers that callers poll to completion,
and OCR, layout and formula-extraction failures remain visible rather than
being stored as empty documents.

An ingest follows a fixed sequence of observable stages: submit, parse, quality
verdict, duplicate decision, split, storage, indexing and a retrieval probe.
Each stage leaves a record that the next stage reads and the audit can revisit:
a parse task identifier, a quality report, a content fingerprint, a split
manifest, a storage record and a probe result. A document counts as ingested
only when a content-derived query retrieves it; an HTTP success from the save
call is not treated as success. The same path accepts scanned pages, Office
files and plain text, and a failure at one stage never silently promotes the
document to the next.

\subsection{Quality gates and ingest lifecycle}

Before storage, duplicate detection checks exact content fingerprints and
flags near-duplicates for review at a 0.90 similarity threshold. After parsing,
a quality gate checks extracted text for OCR noise, binary residue, bodyless
headings, whitespace collapse and language mismatch. The parsed result is saved
with task provenance but is not indexed by that save; direct-text documents use
the create path, which triggers background indexing. The ingestion agent
selects a library whose scope matches the parsed content's subject, not its
filename.

Documents exceeding the 30,000-character reading-context ceiling are split
at structural boundaries. The splitter targets 0.85 utilization and avoids
cuts that would divide atomic units such as equations or tables. Descriptions
and tags are derived from the body and checked against it; mismatches require
metadata correction, not document rewriting. Indexing populates vector, lexical
and graph representations separately; post-index verification checks for a
non-empty vector collection, expected chunk identifiers and graph/document
state. Ingestion also runs content-derived probes before marking an ingest
complete; API success alone is not proof of retrievability.

\subsection{Acquisition evidence and accounting limits}

The 100-paper campaign produced 322 stored parts in five libraries.
Its reconciliation is incomplete: the state file records 91 submitted names;
at the 2026-10-13 evaluation snapshot 79 papers were retrievable and 21 were
not, and
the surviving artifacts do not enumerate attempts and failures per paper, so
the missing 21 cannot be attributed to parsing alone. Formula-dense scans and
a recovered batch-process crash were observed; neither establishes the cause
for every missing paper.

A separate cross-lingual test processed 400 Wikipedia pages per language:
ingestion success 1.0 in both rounds, membership accuracy 0.875 for English
against 1.0 for Chinese and Japanese, with the per-language storage and
self-retrieval results in Section~\ref{sec:e1}. These measurements
characterize ingestion and storage, not answer quality.
